\section*{Chapitre \ref{ch:b2:crustal-melting} --- Fusion de la croûte~: migmatites et granites}

\begin{solution}{exo:b2:crustal-melting:1}
(a) Pour \qty{100}{\gram}~: $n(\ce{Al2O3}) = 17.03/101.96 = \qty{0.16703}{mol}$,
$n(\ce{CaO}) = 0.09183$, $n(\ce{Na2O}) = 0.06783$, $n(\ce{K2O}) = 0.03076$.
A/CNK $= 0.16703/0.19042 = 0.877$~; A/NK $= 0.16703/0.09859 = 1.694$.
(b) A/CNK $< 1 <$ A/NK~: \omterm{def:b2:crustal-melting:asi}{métalumineux}. Le calcium que les feldspaths ne peuvent pas
prendre passe dans des silicates de calcium~: hornblende (ou clinopyroxène, titanite).
\end{solution}

\begin{solution}{exo:b2:crustal-melting:2}
(a) Muscovite et grenat portent de l'aluminium au-delà des feldspaths~: \omterm{def:b2:crustal-melting:asi}{hyperalumineux}~; ce granite
peut être de type S. (b) Hornblende et titanite portent du calcium au-delà des feldspaths~:
\omterm{def:b2:crustal-melting:asi}{métalumineux}. (c) Égirine et une amphibole sodique portent du sodium sans l'aluminium
des feldspaths~: \omterm{def:b2:crustal-melting:asi}{hyperalcalin}.
\end{solution}

\begin{solution}{exo:b2:crustal-melting:3}
Un kilomètre de roche pèse $2800\times9.81\times1000 = \qty{2.747e7}{Pa}$, donc
$P = 0.2747\,z$ (\unit{kbar}, $z$ en \unit{km}). La droite du solidus est $T = 680 - 7.5(P -
2) = 695 - 2.060\,z$. En posant $10 + 20z = 695 - 2.060z$, on trouve $z = 685/22.060 =
\qty{31.1}{km}$, $T = \qty{631}{\celsius}$, $P = \qty{8.5}{kbar}$.
\end{solution}

\begin{solution}{exo:b2:crustal-melting:4}
(a) Les couches blanches sont des \omterm{def:b2:crustal-melting:migmatite}{leucosomes}, les liserés noirs des \omterm{def:b2:crustal-melting:migmatite}{mélanosomes}, le gneiss ordinaire le
mésosome. (b) \omterm{def:b2:crustal-melting:migmatite}{Leucosomes} discrets et bien délimités, \omterm{def:b2:crustal-melting:migmatite}{mélanosomes} et mésosomes~: une \omterm{def:b2:crustal-melting:migmatite}{métatexite},
à structure stromatique. (c) La roche est devenue une \omterm{def:b2:crustal-melting:migmatite}{diatexite} (schlierique et nébulitique)~:
la fusion est allée plus loin, les minéraux sombres ont fondu aussi et le liquide est resté en place,
effaçant la stratification.
\end{solution}

\begin{solution}{exo:b2:crustal-melting:5}
(a) $M = 39.10 + 3\times24.31 + 26.98 + 3\times28.09 + 12\times16.00 + 2\times1.008 =
\qty{417.30}{\gram\per\mole}$~; $w = 18.02/417.30 = \qty{4.32}{\percent}$.
(b) $F = 0.25\times0.0432/0.040 = 0.270$, soit \qty{27}{\percent} de liquide en masse, dans la fourchette des
20 à \qty{30}{\percent} en volume produits par les expériences de \omterm{def:b2:crustal-melting:reactions}{fusion par déshydratation} de la biotite sur des
roches voisines.
\end{solution}

\begin{solution}{exo:b2:crustal-melting:6}
(a) Feldspath potassique \qty{278.35}{\gram\per\mole}, albite 262.24, anorthite 278.22, muscovite
398.33. (b) Quantités~: feldspath potassique $30/278.35 = \qty{0.10778}{mol}$, albite $0.13347$,
anorthite $0.01797$, muscovite $0.01255$. Le feldspath potassique est $\tfrac12$\ce{K2O.}$\tfrac12$\ce{Al2O3},
albite $\tfrac12$\ce{Na2O.}$\tfrac12$\ce{Al2O3}, anorthite \ce{CaO.Al2O3}, muscovite
$\tfrac12$\ce{K2O.}$\tfrac32$\ce{Al2O3} par unité formulaire. $n(\ce{Al2O3}) = 0.05389 +
0.06673 + 0.01797 + 0.01883 = \qty{0.15742}{mol}$~; $n(\ce{K2O}) = 0.05389 + 0.00628 =
0.06017$, $n(\ce{Na2O}) = 0.06673$, $n(\ce{CaO}) = 0.01797$, somme $0.14487$.
A/CNK $= 0.15742/0.14487 = 1.087$. (c) Sans la muscovite, les trois feldspaths donnent
exactement 1 (\cref{prop:b2:crustal-melting:feldspar})~; la muscovite, avec un \ce{Al2O3} de plus que le
feldspath potassique par unité formulaire, rend la roche \omterm{def:b2:crustal-melting:asi}{hyperalumineuse}.
\end{solution}

\begin{solution}{exo:b2:crustal-melting:7}
Basalte~: $900 + 2.5\times(1200 - 850) = \qty{1775}{\mega\joule\per\cubic\metre}$. Croûte~:
$2.5\times(850 - 650) + 0.20\times900 = \qty{680}{\mega\joule\per\cubic\metre}$.
$V_c/V_b = 1775/680 = 2.6$. L'estimation ignore la chaleur qui fuit par conduction
vers le haut et latéralement, ainsi que le temps que met la chaleur à se répandre~; elle suppose aussi que toute
la chaleur du basalte atteint une croûte prête à fondre~: c'est une borne supérieure.
\end{solution}

\begin{solution}{exo:b2:crustal-melting:8}
(a) $n(\ce{Al2O3}) = 14.15/101.96 = 0.13878$, $n(\ce{CaO}) = 0.01753$, $n(\ce{Na2O}) =
0.07809$, $n(\ce{K2O}) = 0.04416$~; A/CNK $= 0.13878/0.13978 = 0.993$, A/NK $=
0.13878/0.12225 = 1.135$~: \omterm{def:b2:crustal-melting:asi}{métalumineux}, de peu. (b) Avec \qty{14.30}{\percent} \ce{Al2O3},
$0.14025/0.13978 = 1.003$~; avec \qty{14.00}{\percent}, $0.13731/0.13978 = 0.982$. L'indice
chevauche 1 dans l'incertitude certifiée~: l'obsidienne se trouve à la frontière métalumineux--hyperalumineux,
et sa classe n'est pas robuste.
\end{solution}

\begin{solution}{exo:b2:crustal-melting:9}
(a) Moyennes $\bar P = \qty{5.0}{kbar}$, $\bar w = \qty{10.045}{\percent}$~;
$\sum(P - \bar P)(w - \bar w) = 17.62$, $\sum(P - \bar P)^2 = 14.0$~; pente
$17.62/14.0 = \qty{1.259}{\percent}$ par kilobar, ordonnée à l'origine $10.045 - 1.259\times5.0 =
\qty{3.75}{\percent}$. La pente concorde avec le 1.27 cité, à
\qty{1}{\percent} près. (b) $3.75 + 1.259\times6 = \qty{11.30}{\percent}$. (c) $10/0.1130 =
\qty{88.5}{\gram}$ de liquide. (d) La droite repose sur des données entre 3 et \qty{8}{kbar} à
\qty{800}{\celsius} seulement~; rien ne garantit que la solubilité reste linéaire à plus haute
pression, où le liquide et le fluide se ressemblent davantage.
\end{solution}

\begin{solution}{exo:b2:crustal-melting:10}
(a) La droite monte de \qty{4}{kbar} sur \qty{50}{\celsius}~: $P = 6 + 0.08\,(T - 750)$~;
à \qty{770}{\celsius}, $P = \qty{7.6}{kbar}$. (b) De 7.6 à \qty{6}{kbar}, la température descend le long de la droite de 770 à
\qty{750}{\celsius}~; la chaleur dégagée, $2.5\times20 = \qty{50}{\mega\joule\per\cubic\metre}$,
fait fondre $50/900 = 0.056$, soit environ
\qty{6}{\percent} de la roche en volume. (c) C'est la fenêtre proposée pour les
\omterm{def:b2:crustal-melting:leucogranite}{leucogranites} de l'Himalaya, 6 à \qty{8}{kbar} et 750 à \qty{770}{\celsius}~: la décompression
seule produit du liquide, mais seulement quelques pour cent à chaque franchissement de la courbe.
\end{solution}

\begin{solution}{exo:b2:crustal-melting:11}
$\ell^2/\kappa = (2000)^2/10^{-6} = \qty{4e12}{s} = \qty{0.13}{Ma}$ pour le \omterm{def:b2:magma-emplacement:bodies}{sill}~;
$(10^4)^2/10^{-6} = \qty{1e14}{s} = \qty{3.2}{Ma}$ pour le sous-placage. Ces deux durées sont courtes
à côté de dix Ma~: un seul événement ne peut pas maintenir la croûte en fusion aussi longtemps, ce qui exige une chaleur
apportée à plusieurs reprises, par des intrusions répétées ou par le réchauffement lent d'une
croûte épaissie.
\end{solution}

\begin{solution}{exo:b2:crustal-melting:12}
(a) $5.0x + 2.0(1 - x) = 2.75$ donne $x = 0.75/3.0 = 0.25$~: un quart du néosome est
\omterm{def:b2:crustal-melting:migmatite}{leucosome}. (b) $0.25\times0.20 + 0.75\times4.0 = 0.05 + 3.00 = \qty{3.05}{\percent}$, comme
dans le mésosome~: l'hypothèse de système clos est confirmée. (c) Pour \qty{100}{\gram} de
roche mère, \qty{25}{\gram} de liquide et \qty{75}{\gram} de résidu se sont formés~; si $m$ grammes de
liquide sont sortis, $(25 - m)/(100 - m) = 0.10$, donc $25 - m = 10 - 0.1m$ et $m = 15/0.9 =
\qty{16.7}{\gram}$~: deux tiers du liquide ont quitté l'affleurement.
\end{solution}

\begin{solution}{pb:b2:crustal-melting:1}
\noindent\textbf{Partie I.}

\textbf{1.} $39.10 + 3\times26.98 + 3\times28.09 + 12\times16.00 + 2\times1.008 =
\qty{398.33}{\gram\per\mole}$.

\textbf{2.} $18.02/398.33 = \qty{4.52}{\percent}$.

\textbf{3.} $200\times0.0452 = \qty{9.04}{\gram}$ d'eau par kilogramme de schiste.

\textbf{4.} $9.97 + (13.82 - 9.97)\times2/3 = \qty{12.54}{\percent}$. Le liquide minimal vaut
$9.04/0.1254 = \qty{72.1}{\gram}$ par kilogramme~: $F \ge \qty{7.21}{\percent}$
(\cref{prop:b2:crustal-melting:water}).

\textbf{5.} $9.04/0.060 = \qty{151}{\gram}$~: $F = \qty{15.1}{\percent}$.

\textbf{6.} Par exemple~: la roche peut contenir moins de muscovite~; une partie de la muscovite peut
déjà s'être décomposée sous le solidus, avec le quartz, en feldspath potassique et alumino-silicate,
son eau quittant la roche sous forme de fluide avant le début de la fusion~; le liquide peut contenir plus d'eau
que \qty{6}{\percent}~; une partie de l'eau peut rester dans la biotite.

\noindent\textbf{Partie II.}

\textbf{7.} $680 - 7.5\times(7.8 - 2) = \qty{636.5}{\celsius}$~; la roche culmine
à \qty{43.5}{\celsius} en dessous et ne peut pas fondre, même en présence d'eau libre.

\textbf{8.} $605 + (19.2 - 15.3)/(27 - 15.3)\times65 = \qty{626.7}{\celsius}$. À
$t = 0$, la base, à \qty{595}{\celsius}, est \qty{32}{\celsius} en dessous~; à
\qty{10}{Ma}, à \qty{656}{\celsius}, elle est au-dessus.

\textbf{9.} La fusion au solidus humide exige de l'eau libre, et une roche métamorphique chauffée
en contient très peu~: l'eau libre est consommée par le premier liquide, qui cesse alors de croître
jusqu'à ce que les micas se décomposent.

\textbf{10.} Entre 30 et \qty{40}{Ma}~: $30 + (800 - 795)/(845 - 795)\times10 =
\qty{31.0}{Ma}$.

\textbf{11.} La base passe sous \qty{12.5}{kbar} entre 70 et \qty{80}{Ma}~:
$70 + (13.7 - 12.5)/(13.7 - 12.4)\times10 = \qty{79.2}{Ma}$, à
$914 - 14\times0.923 = \qty{901}{\celsius}$, dans l'encadrement de l'amphibole~: une amphibolite
située à la base serait en train de fondre.

\textbf{12.} $90 + (11.0 - 10)/(11.0 - 9.6)\times10 = \qty{97.1}{Ma}$, à $867 -
53\times0.714 = \qty{829}{\celsius}$. C'est au-dessus de l'encadrement de la muscovite~:
la muscovite s'est décomposée. Cela se situe dans l'intervalle de 800 à \qty{850}{\celsius} de la décomposition de la biotite
(mesurée à plus basse pression)~: la biotite, l'autre minéral hydraté du schiste, est à sa
limite.

\noindent\textbf{Partie III.}

\textbf{13.} $n(\ce{Al2O3}) = 15.0/101.96 = \qty{0.14712}{mol}$, $n(\ce{CaO}) =
0.55/56.08 = 0.00981$, $n(\ce{Na2O}) = 3.4/61.98 = 0.05486$, $n(\ce{K2O}) = 4.9/94.20 =
0.05202$.

\textbf{14.} A/CNK $= 0.14712/0.11669 = 1.261$.

\textbf{15.} A/NK $= 0.14712/0.10688 = 1.377$~; A/CNK $> 1$~: \omterm{def:b2:crustal-melting:asi}{hyperalumineux}.

\textbf{16.} \omterm{def:b2:crustal-melting:asi}{Hyperalumineux}, issu d'un schiste~: un \omterm{def:b2:crustal-melting:types}{granite de type S}. Son excès d'aluminium se trouve dans
la muscovite et le grenat (ou la cordiérite).

\textbf{17.} $0.14712 - 0.11669 = \qty{0.03043}{mol}$ de \ce{Al2O3}~; en muscovite,
$0.03043\times398.33 = \qty{12.1}{\gram}$ dans \qty{100}{\gram} de granite.

\textbf{18.} L'obsidienne est (de peu) \omterm{def:b2:crustal-melting:asi}{métalumineuse}, sans l'excès d'aluminium d'un liquide
métasédimentaire~: sa source est probablement différente. Mais un seul échantillon prouve peu~:
toute suite de granitoïdes compte des membres légèrement \omterm{def:b2:crustal-melting:asi}{hyperalumineux}, et une suite se juge sur
l'ensemble de sa tendance.

\noindent\textbf{Partie IV.}

\textbf{19.} $10^{12}\,\unit{\cubic\metre}\times2650 = \qty{2.65e15}{kg}$.

\textbf{20.} Le granite est le liquide moins son eau~: $2.65\times10^{15}/0.94 =
\qty{2.82e15}{kg}$ de liquide.

\textbf{21.} Seule une fraction $0.151\times0.70$ de la masse du schiste part en liquide~:
$2.82\times10^{15}/(0.151\times0.70) = \qty{2.67e16}{kg}$.

\textbf{22.} $2.67\times10^{16}/2800 = \qty{9.5e12}{\cubic\metre}$, soit
\qty{9.5e3}{\cubic\kilo\metre}, près de dix fois le \omterm{def:b2:magma-emplacement:bodies}{pluton}.

\textbf{23.} $9.5\times10^3/5 = \qty{1.9e3}{\square\kilo\metre}$, soit un carré d'environ
\qty{44}{km} de côté.

\textbf{24.} Avec $F = \qty{7.21}{\percent}$, la source est $15.1/7.21 = 2.1$ fois
plus grande, soit environ \qty{2.0e4}{\cubic\kilo\metre}. Le volume de croûte source fondu pour faire
le \omterm{def:b2:magma-emplacement:bodies}{pluton}, avec le \omterm{def:b2:mantle-melting:fraction}{taux de fusion} de la question~5, est
\textbf{\qty[text-series-to-math=true]{9.5e3}{\cubic\kilo\metre}}.
\end{solution}
